\documentclass{report}
\usepackage{amsmath,amssymb}
\setlength{\parindent}{0mm}
\setlength{\parskip}{1em}
\begin{document}
\begin{center}
\rule{6in}{1pt} \
{\large Carmen Rodrigo \\
{\bf An efficient multigrid algorithm for finite difference discretizations on triangular grids}}

Department of Applied Mathematics \\ University of Zaragoza \\ C/ Maria de Luna \\ 3 \\ 50018 Zaragoza \\ Spain
\\
{\tt carmenr@unizar.es}\\
Francisco J. Gaspar\\
Francisco J.  Lisbona\end{center}

The differential operators divergence, gradient and rotor are often used
to formulate mathematical physics problems. To approximate the solution
of this type of problems it is convenient to use irregular grids to fit
better the spatial domain. A natural way to discretize the differential
problem is to define the corresponding discrete grid operators. These
must satisfy the main properties of the continuous operators and also
some compatibility
relations between them. In the literature the associated methods are
called mimetic finite difference methods. In a recent paper [1], a
specially clear approach, called VAGO (Vector Analysis Grid Operators)
method, has been proposed on Delaunay triangular grids. For this
discretization, it is not necessary to use a concrete coordinate system.
This fact is specially interesting when computational irregular grids are
considered.

An important issue in solving this type of problems is the efficient
solution of the corresponding algebraic system of equations which results
after the discretization process, and multigrid methods are generally
accepted as the fastest numerical methods for the solution of elliptic
partial differential equations. Although the algebraic multigrid is a
useful tool to solve problems on unstructured grids, we consider in this
work an efficient and robust geometric multigrid algorithm to solve
problems on triangular grids in a free-matrix version. This is achieved
by using an adequate data structure resulting in a very fast solver with
very low memory requirements. Finally, we consider some test problems to
show the goodness of this multigrid method.

[1] P.N. Vabishchevich, Finite difference approximation of mathematical
physics problems on irregular grids,
Comp. Meth. Appl. Math. 5 (2005) pp. 294--330.


\end{document}
